The large integration of microphones into devices increases the opportunities for Acoustic Side-Channel Attacks (ASCAs), as these can be used to capture keystrokes' audio signals that might reveal sensitive information. However, the current State-Of-The-Art (SOTA) models for ASCAs, including Convolutional Neural Networks (CNNs) and hybrid models, such as CoAtNet, still exhibit limited robustness under realistic noisy conditions. Solving this problem requires either: (i) an increased model's capacity to infer contextual information from longer sequences, allowing the model to learn that an initially noisily typed word is the same as a futurely collected non-noisy word, or (ii) an approach to fix misidentified information from the contexts, as one does not type random words, but the ones that best fit the conversation context. In this paper, we demonstrate that both strategies are viable and complementary solutions for making ASCAs practical. 
We observed that no existing solution leverages advanced transformer architectures' power for these tasks and propose that: (i) Visual Transformers (VTs) are the candidate solutions for capturing long-term contextual information and (ii) transformer-powered Large Language Models (LLMs) are the candidate solutions to fix the ``typos'' (mispredictions) the model might make. Thus, we here present the first-of-its-kind approach that integrates VTs and LLMs for ASCAs.

We first show that VTs achieve SOTA performance in classifying keystrokes when compared to the previous CNN benchmark. 
Second, we demonstrate that LLMs can mitigate the impact of real-world noise. Evaluations on the natural sentences revealed that: (i) incorporating LLMs (e.g., GPT-4o) in our ASCA pipeline boosts the performance of error-correction tasks; and (ii) the comparable performance can be attained by a lightweight, fine-tuned smaller LLM (67 times smaller than GPT-4o), using Low-Rank Adaptation (LoRA). Our results and findings highlight the practical viability of our solution toward effective ASCA.